\renewcommand{\thechapterdisplay}{Vigésimo quarto dia}
\renewcommand{\thechapterrunner}{Quinto modo de aliviar as almas do Purgatório: o sofrimento}
\chapter[Vigésimo quarto dia --- Quinto modo de aliviar as almas do Purgatório: o sofrimento]{Quinto modo de aliviar as almas do Purgatório: o sofrimento}
\section*{Sofrimento voluntário}

Aliviemos as almas do Purgatório, diz São João Crisóstomo, aliviemo-las por tudo o que sofremos, pois Deus se encarrega de aplicar aos mortos os méritos dos vivos. O sofrimento é a grande satisfação que o Senhor pede às almas cujo amor é devedor de sua justiça. Portanto, soframos por elas afim de que sofram menos. Ó, se tivéssemos uma fé mais firme, uma caridade mais viva, que mortificações não nos imporíamos para aliviar e libertar os familiares, os amigos que tanto nos amaram, e que sofrem agora de uma maneira tão horrível? A penitência, o jejum, as austeridades, seriam nossas práticas corriqueiras. Tenhamos ao menos a coragem de realizar alguns sacrifícios leves: privar-se de um prazer, evitar uma afeição perigosa, cessar más leituras, maus hábitos, renunciar objetos de luxo e de pura vaidade. Dizia o padre Félix: \textit{Escolhei a melhor vítima! Escolhei-a sobretudo no fundo de vossos corações. Por aqueles que vós mais amais, sacrificai o que tendes de mais caro; sacrificai-vos a vós mesmos, e que o preço do sacrifício pessoal possa se tornar o fim do sofrimento de nossos queridos falecidos.}

Ah! Eu as vejo --- essas almas bem-aventuradas --- eu as vejo subir para o Céu sobre as asas de nossos sacrifícios, de nossas austeridades, de nossos sofrimentos. Elas voam triunfantes e nos agradecem pela nossa generosidade. E quando estiverem na glória, elas nos devolverão abundantemente o que tivermos feito por elas. Que motivo de consolo e de esperança! Ó meu Deus! Ó Jesus crucificado! Fazei-nos compreender o valor do sofrimento.


\section*{Sofrimento involuntário}

Se nos falta coragem para o sofrimento voluntário, a Providência Divina nos impõe sofrimentos mais meritórios para nós e para os mortos, porque não são escolhidos por nós. São as aflições, as dores do espírito, do coração e do corpo, inevitáveis neste mundo. Infelizmente, as encontramos em todo lugar, em todos os estados de vida, em todas as condições. Nossa vida na terra é um combate diário, um longo e sofrido martírio. Devemos nós reclamar? Não, pois todas as nossas dores podem se tornar um meio de salvação para nós e para os outros; e podemos nos servir delas para diminuir as mais cruéis de todas as dores: aquelas a que são submetidas as almas do Purgatório. Sim, com as cruzes que a Providência coloca em nossos ombros, com os espinhos que cravam nossos corações, com uma lágrima, com um suspiro, com um ato de resignação, nos podemos amenizar as grandes misérias do além-túmulo e secar as lágrimas de nossos conhecidos que passaram. 

Coragem, pois, alma cristã, suportemos um pouco de frio, e refrescaremos as vítimas que queimam no meio do fogo da cólera divina. Soframos um pouco de calor, e trocaremos os ardores desse fogo em um suave orvalho. Suportemos um incômodo, e arrancaremos almas do mais profundo abismo. Aceitemos um cansaço, e as levaremos sobre tronos de glória no Céu. Para nós, um momento de dor; para elas, uma eternidade de felicidade!


\section*{Exemplo }

Um doente, conta Santo Antonino, estava sofrendo terrivelmente e pedia a Deus, com lágrimas, a cura de seus males. Um anjo lhe apareceu e disse: \textit{O Senhor me enviou a vós, para dar-lhe a escolha entre um ano de sofrimento na terra, ou um dia no Purgatório.} O doente não hesitou: \textit{Um dia no Purgatório! Pelo menos terminarão minhas dores.} Ele expirou imediatamente, e sua alma foi precipitada no abismo de expiação. Então o anjo compassivo foi até ele para o consolar. Quando viu o anjo, o miserável solta um grito lancinante, semelhante ao rugido do inferno, e exclama: \textit{Anjo sedutor, vós me enganastes! Vós me assegurastes que eu não ficaria mais que um dia no purgatório, e eis que há mais de vinte anos que eu estou aqui nos mais horríveis suplícios!} --- \textit{Pobre infeliz}, respondeu o anjo, \textit{está errado. O rigor de seus tormentos o faz exagerar a duração, e sentir como um século o que não passa de um instante. Não se engane, apenas alguns minutos se passaram desde que morreu, e seu cadáver ainda não esfriou no leito de morte.} --- \textit{Então, consiga que eu retorne à terra para ali sofrer, durante um ano, tudo o que Deus quiser.} Seu pedido foi atendido. O doente exortava todos os que vinham vê-lo a aceitar de bom coração todas as penas deste mundo, em vez de se exporem aos tormentos do outro. \textit{A paciência nas dores}, dizia ele frequentemente, \textit{é a chave de ouro do Paraíso.} Aproveitemos, portanto, para nós mesmos e para as almas do Purgatório, os sofrimentos que a Providência julgar apropriado nos enviar.


\textbf{Oremos.} Sejais bendito, ó meu Deus, que desejastes que os sofrimentos e as dores incessantes de minha vida se tornem para mim uma abundantes fonte de méritos, e um meio de satisfazer à vossa justiça pelas almas que me são queridas. De agora em diante, em vez de me queixar do peso das minhas cruzes, eu as suportarei com paciência e resignação, e vós lançareis sobre mim e sobre meus parentes falecidos um olhar de misericórdia. Ó Jesus, sede-lhes propício! Chamai para perto de Vós os vossos filhos, nossos irmãos! Que descansem em paz!



\begin{center}
\textit{Requiescant in pace!}
\end{center}

Rezai agora uma dezena do terço e as orações no final deste livro.
